\section{Computing \texorpdfstring{$A(p)$}{A(p)} in practice}
\label{a:sec:practical}

In the absence of an exact formula we are left with approximation, and the two
regimes identified above suggest how to split the work.

Over most of the interval, direct iteration is entirely satisfactory: the infinite
product \cref{a:eq:prodFormula} converges quickly, its partial products bound $A(p)$
from above at every order, the tail bound \cref{a:eq:producttailbound} supplies a
matching lower bound and hence a stopping certificate, and modest numbers of
iterations give many digits. The difficulty is confined to a neighbourhood of
$p=\tfrac12$, and it has two parts. First, as $A(p)\to0$ the quantity of interest
$1/A$ diverges, so a fixed absolute error in $A$ becomes an unbounded error in the
quasi-stationary mean. Second, and this is the binding constraint, as the process is
defined ever closer to criticality its lifetime, though still finite, grows
enormously, and $S_n$ recedes to zero ever more slowly. The iteration must be
carried to a generation at which $S_n$ is essentially unchanging, and
\cref{a:tab:Avals} shows what that costs: about $10^2$ iterations at $p=0.2$,
$10^3$ at $p=0.45$, but nearly $5\times10^5$ at $p=0.4999$.

The remedy is to hand the near-critical region to the asymptotic
\cref{a:eq:Anearcritical}, which is exact in the limit and cheap everywhere:
\begin{equation}
\label{a:eq:Apractical}
A(p) \approx
\begin{cases}
\displaystyle \frac12\prod_{k=1}^{N}\lrb{1-\frac{S_k}{2}}, & p<p_\ast,\quad N\gg1,\\[14pt]
2-4p, & p\ge p_\ast,
\end{cases}
\end{equation}
with the crossover $p_\ast$ chosen just below $\tfrac12$. \Cref{a:rem:rate}
indicates where to put it. The relative error of the asymptotic is about
$\varepsilon\log(1/\varepsilon)$ by \cref{a:thm:nearcritical}, so at
$p_\ast = 0.4999$ (that is $\varepsilon = 2\times10^{-4}$) the asymptotic is already
accurate to about two parts in $10^3$, while the product still converges in a few
hundred thousand iterations. Any $p_\ast$ in the range $0.499$ to $0.4999$ gives a
scheme accurate to better than one per cent throughout, and the two branches can be
blended over a narrow window if continuity of the derivative matters.

Three practical points are worth recording.

First, use the product rather than the raw ratio $S_n/(2p)^n$. Both require the same
iteration, but the product is monotone decreasing in $N$ and comes with the two-sided
certificate \cref{a:eq:producttailbound}, whereas the ratio approaches its limit
from no fixed side and offers no error control.

Second, stop on the certificate, not on apparent convergence. In the near-critical
regime successive floating-point iterates cease to differ long before the tail is
genuinely negligible, and a stopping rule based on observed change silently returns
the wrong answer. The bracket in \cref{a:eq:producttailbound} costs one extra
multiplication and removes the failure mode.

Third, if the quantity actually needed is the quasi-stationary mean $1/A$ rather
than $A$ itself --- as it usually is --- then the series \cref{a:eq:Aseries}
computes it directly, and its partial sums increase to the answer, giving the
complementary one-sided bound.

This is the recipe used for the numerical values quoted throughout this chapter
and in \ChM, including \cref{a:tab:Avals}.
